
\documentclass[11pt]{article}
\usepackage[margin=1in]{geometry}
\usepackage{amsmath,amssymb,booktabs,graphicx,xcolor,caption,enumitem}
\usepackage[colorlinks=true,linkcolor=blue,citecolor=blue,urlcolor=blue]{hyperref}
\usepackage{listings}
\definecolor{codekw}{rgb}{0.15,0.15,0.6}
\definecolor{codecm}{rgb}{0.30,0.55,0.30}
\definecolor{codestr}{rgb}{0.6,0.2,0.2}
\definecolor{codebg}{rgb}{0.97,0.97,0.97}
\definecolor{darkgreen}{rgb}{0.0,0.5,0.0}
\lstset{
  basicstyle=\ttfamily\small, keywordstyle=\color{codekw}\bfseries,
  commentstyle=\color{codecm}\itshape, stringstyle=\color{codestr},
  backgroundcolor=\color{codebg}, frame=single, framerule=0.2pt, rulecolor=\color{gray},
  breaklines=true, showstringspaces=false, columns=fullflexible, language=Python,
  numbers=left, numberstyle=\tiny\color{gray}, xleftmargin=2.2em,
}
\newcommand{\F}{\mathbb{F}}
\newcommand{\Sp}{\operatorname{Sp}}
\newcommand{\code}[1]{\texttt{#1}}
\newcommand{\file}[1]{\texttt{#1}}
\captionsetup{font=small}

\title{\textbf{SparseGF2: Circuit Construction}\\[2pt]
\large An Implementation Reference}
\author{MIPT--QEC project}
\date{\today}

\begin{document}
\maketitle

\begin{abstract}
\noindent This document specifies, in complete implementation detail, how the SparseGF2 stabilizer simulator
\emph{constructs and executes} the monitored random-Clifford circuits used in the measurement-induced phase
transition (MIPT) study. It is the software companion to the mathematical specification
(\file{docs/sparsegf2\_mathematics.tex}): where that document derives the stabilizer/symplectic formalism, this
one traces the actual code path --- configuration, graph and edge selection, the sparse tableau data structure,
the purification picture and initial scramble, the layer-by-layer execution loop, gate application through the
Clifford lookup table and Numba kernels, single-qubit measurement, observables/order parameters, and the record
and reproducibility conventions. Every component is tied to the concrete module, class, and function that
implements it. File references are relative to \file{src/sparsegf2/}.
\end{abstract}

\tableofcontents
\clearpage

% ======================================================================
\section{Overview: from specification to executed circuit}
\label{sec:overview}
A run is a pipeline of stages, each owned by one module. A circuit is \emph{never} materialized as a static list
of gates; it is generated on the fly, layer by layer, directly against the evolving stabilizer state.
\begin{enumerate}[itemsep=2pt]
\item \textbf{Configuration} (\file{circuits/config.py}, \code{CircuitConfig}) --- an immutable, validated
specification: the number of system qubits $n$, the depth mode and factor, the measurement rate $p$, the
interaction-graph family, the gating mode, the measurement-eligibility mode, the ``picture'' (plain vs.\
purification), the scramble/warmup options, and the seed.
\item \textbf{Interaction graph} (\file{circuits/graphs.py}, \code{GraphTopology}) --- the set of qubit pairs on
which two-qubit gates may act, plus (for matching-based gating) an edge coloring / $1$-factorization.
\item \textbf{Scheduler} (\file{circuits/scheduler.py}, \file{circuits/matching.py}) --- turns the graph into a
per-layer sequence of gated edges (with a Clifford index each) and a per-layer list of measurement candidates,
fixing the RNG consumption order.
\item \textbf{Stabilizer state} (\file{core/sparse\_tableau.py}, \code{SparseGF2}) --- the Aaronson--Gottesman
tableau stored sparsely over $\F_2$, with an optional dense/hybrid bit-packed mirror.
\item \textbf{Picture / initial state} (\file{circuits/picture.py}, \code{setup\_picture}) --- builds the initial
state (e.g.\ $n$ Bell pairs across a system/reference split for the purification picture) and reports the layout
and order parameter.
\item \textbf{Runner} (\file{circuits/runner.py}, \code{SimulationRunner}) --- the execution loop: optional
scramble, warmup, then for each measured layer apply the gate step and the measurement step, evaluate the
order parameter, take requested checkpoints, and support early stopping.
\item \textbf{Records} (\file{circuits/records.py}, \code{SampleRecord}) --- the per-run outputs.
\end{enumerate}
Section~\ref{sec:changes} separates the stable simulator semantics from the additive circuit and checkpoint
extensions used by the deep purification-time campaign.

% ======================================================================
\section{Run configuration}
\label{sec:config}
Everything about a run is fixed by one immutable dataclass, \code{CircuitConfig} (\file{config.py}), validated in
\code{\_\_post\_init\_\_}. Table~\ref{tab:config} lists the fields.

\begin{table}[h]\centering\small
\begin{tabular}{lll}
\toprule
Field & Default & Role\\
\midrule
\code{graph\_spec} & --- & Graph of allowed 2-qubit edges (\code{str} / \code{GraphTopology} / \code{nx.Graph})\\
\code{n} & --- & Number of system qubits\\
\code{picture} & \code{pure\_state} & \code{pure\_state} / \code{purification} / \code{single\_ref} (Sec.~\ref{sec:picture})\\
\code{gating\_mode} & \code{brickwork} & \code{brickwork} / \code{random\_edge} / \code{random\_pool} / \code{all\_edges}\\
\code{matching\_mode} & \code{round\_robin} & brickwork only: \code{round\_robin} / \code{palette} / \code{fresh}\\
\code{gates\_per\_layer} & \code{None} & random-edge modes: default $1$ (\code{random\_edge}) or $n/2$ (\code{random\_pool})\\
\code{measurement\_mode} & \code{bernoulli} & candidates: \code{bernoulli} (all) / \code{gated} / \code{random\_pair} / \code{uniform\_count}\\
\code{meas\_count} & \code{None} & for \code{uniform\_count}; default $n/2$\\
\code{p} & \code{0.15} & per-candidate measurement probability $\in[0,1]$\\
\code{depth\_mode} & \code{O(n)} & \code{O(n)} / \code{O(log\_n)} / \code{until\_purified} (adaptive; needs a reference)\\
\code{depth\_factor} & \code{8} & multiplier in the depth formula ($\ge1$)\\
\code{total\_layers\_override} & \code{None} & optional literal positive measured-layer count; bypasses the depth formula\\
\code{n\_cliffords} & \code{720} & size of the sampled 2-qubit Clifford set, $\in[1,720]$\\
\code{base\_seed} & \code{42} & nonnegative base component of every pair-keyed per-sample RNG stream\\
\code{warmup\_layers} & \code{0} & gate-only pre-scrambling layers before the measured phase\\
\code{scramble} & \code{False} & apply one global random Clifford to the system qubits pre-circuit\\
\code{scramble\_entangled\_qubit} & \code{True} & \code{single\_ref}: include the reference's partner in the scramble\\
\code{record\_time\_series} & \code{False} & record the order parameter after every layer\\
\code{pivot\_mode} & \code{None} & measurement pivot: \code{min\_weight} (sparser) / \code{first} (cheaper)\\
\code{use\_numba} & \code{None} & JIT toggle; \code{None} auto-detects\\
\code{hybrid} & \code{False} & enable the sparse/dense bit-packed hybrid path\\
\bottomrule
\end{tabular}
\caption{\code{CircuitConfig} fields.}
\label{tab:config}
\end{table}

\paragraph{Derived quantities.} If \code{total\_layers\_override} is set, \code{total\_layers()} returns that
literal scheduled measured-layer count before considering any depth mode or random-edge rescaling. Otherwise, for
\code{O(n)} it starts from $\mathrm{depth\_factor}\cdot n$, for \code{O(log\_n)} from
$\mathrm{depth\_factor}\cdot\lceil\log_2 n\rceil$, and \code{until\_purified} uses the $O(n)$ value as its cap.
The random-edge modes rescale those formula-based counts to preserve a brickwork-equivalent gate budget;
\code{all\_edges} does not. Adaptive purification or a \code{CHECKPOINT\_STOP} callback may execute fewer layers
than the scheduled cap.
\code{total\_qubits()} is $n$, $2n$, or $n{+}1$ by picture. \code{scramble\_qubits()} returns which qubits the
global scramble covers (\code{range(n)} for every case except \code{single\_ref} with
\code{scramble\_entangled\_qubit=False}, which holds out the reference's partner). The resolved graph is built
once and cached on \code{\_graph} so all samples of a config share it. A stochastic string graph specification
uses \code{base\_seed} alone as its quenched graph-realization seed.

% ======================================================================
\section{The stabilizer state: the sparse tableau}
\label{sec:tableau}
The state is an Aaronson--Gottesman tableau of $2n$ generators ($n$ destabilizers, $n$ stabilizers), phase-free
over $\F_2$ (SparseGF2 tracks only the symplectic $[X|Z]$ part; signs are dropped, which is sufficient for the
entanglement/purification observables). The representation (\code{SparseGF2}, \file{core/sparse\_tableau.py}) is
\emph{sparse}: rather than a dense $2n\times2n$ matrix, it keeps, for each generator, only the qubits on which it
acts, together with inverted indices for $O(1)$ lookups. Table~\ref{tab:tableau} lists the arrays.

\begin{table}[h]\centering\small
\begin{tabular}{lll}
\toprule
Array & Shape & Purpose\\
\midrule
\code{plt} & $(2n,n)$ \code{uint8} & Pauli lookup table: 2-bit code $(x{\ll}1)\,|\,z$, i.e. $I{=}0,Z{=}1,X{=}2,Y{=}3$\\
\code{supp\_q} / \code{supp\_len} & $(2n,n)$/$(2n,)$ & per-generator support (nonzero qubits) and its length\\
\code{supp\_pos} & $(2n,n)$ & back-pointer into \code{supp\_q} for $O(1)$ removal\\
\code{inv} / \code{inv\_len} & $(n,2n)$/$(n,)$ & per-qubit inverted index: which generators touch qubit $q$\\
\code{inv\_pos} & $(2n,n)$ & back-pointer for \code{inv}\\
\code{inv\_x} / \code{inv\_x\_len} & $(n,2n)$/$(n,)$ & per-qubit index of generators with the $x$-bit set on $q$\\
\code{inv\_x\_pos} & $(2n,n)$ & back-pointer for \code{inv\_x}\\
\bottomrule
\end{tabular}
\caption{The seven-array sparse tableau, kept in sync by $O(1)$ inverted-index primitives.}
\label{tab:tableau}
\end{table}

\paragraph{Row layout \& initialization.} Rows $0\ldots n{-}1$ are destabilizers ($X_i$ on qubit $i$); rows
$n\ldots 2n{-}1$ are stabilizers ($Z_i$ on qubit $i$). \code{\_init\_zero\_state} sets \code{plt[i,i]=X} and
\code{plt[n+i,i]=Z} with unit support, which is the $|0^n\rangle$ tableau. The \code{inv\_x} index (generators
that anticommute with a $Z$-measurement on $q$) is what makes measurement cheap (Sec.~\ref{sec:measure}).

\paragraph{Hybrid dense mirror.} With \code{hybrid=True} the tableau additionally maintains a bit-packed
\code{uint64} dense mirror; volume-law (dense) circuits then use packed row operations while measurement-heavy
(sparse) circuits use the sparse path. Correctness is identical; only the constant factors differ.

\paragraph{Constructor.} \code{SparseGF2(n, rng=None, pivot\_mode=None, use\_numba=None, hybrid=False)}. The
\code{rng} is the \emph{measurement-outcome} generator (Sec.~\ref{sec:records}), independent of circuit
construction. Key methods: \code{copy()}, \code{apply\_clifford(S, qubits)} and \code{\_dispatch\_2q(qi,qj,lut)}
(Sec.~\ref{sec:gates}), \code{measure\_z(q)} (Sec.~\ref{sec:measure}), \code{active\_count()} (mean generators
per qubit diagnostic), and \code{to\_symplectic()} $\to (2n,2n)$ \code{uint8} $[X|Z]$.

% ======================================================================
\section{Interaction graph and per-layer edge selection}
\label{sec:graph}
\code{GraphTopology} (\file{graphs.py}) carries \code{name}, \code{n}, a canonical \code{edges} list ($u<v$),
\code{is\_stochastic} (whether a realization depends on the seed, e.g.\ Newman--Watts), an optional
\code{one\_factorization} (partition into perfect matchings, for brickwork), the \code{graph6} string, and an
optional \code{fresh\_matching\_sampler}; \code{chi\_prime} is the number of color classes. Graphs are built
from a spec by \code{from\_spec} (families: \code{cycle}, \code{path}, \code{complete}, \code{lattice\_2d},
\code{newman\_watts}, \dots) or adapted from a \code{networkx} graph by \code{from\_networkx}.

\paragraph{Per-layer selection.} \code{select\_matching(graph, mode, layer\_index, rng)}
(\file{matching.py}) returns a matching: \code{round\_robin} cycles deterministically through the
$1$-factorization (layer $t$ uses class $t\bmod\chi'$, no RNG draw); \code{palette} draws a random class per
layer; \code{fresh} samples a uniformly random perfect matching. The random-edge modes instead draw edge
indices directly. The \code{all\_edges} mode fires every edge in the graph's stored order on every layer; its
placement is deterministic and consumes no RNG.

\paragraph{The scheduler.} \code{CircuitBuilder(config, sample\_seed)} (\file{scheduler.py}) owns one
\code{np.random.default\_rng([base\_seed, sample\_seed])} consumed in a \emph{fixed} per-layer order ---
(1)~gate placement, (2)~one Clifford index per gate, (3)~measurement-candidate draws --- so a run is
bit-reproducible. It yields \code{CircuitLayer} objects (fields: \code{gate\_pairs}, \code{cliff\_indices},
\code{meas\_qubits}, \code{meas\_candidates}) via \code{warmup\_layers\_iter()} (gate-only) and \code{layers()}
(measured).

% ======================================================================
\section{The purification picture and the initial scramble}
\label{sec:picture}
\code{setup\_picture(picture, n\_system, *, rng, pivot\_mode, use\_numba, hybrid)} (\file{picture.py}) returns a
freshly initialized \code{SparseGF2} with the picture's setup operations already applied, and a \code{PictureSpec}
describing the layout (\code{system\_qubits}, \code{reference\_qubits}, \code{total\_qubits}) and the
\code{order\_parameter} name. The three pictures:
\begin{itemize}[itemsep=2pt]
\item \code{pure\_state}: $n$ qubits in $|0^n\rangle$; no setup ops; order parameter is the half-cut entropy.
\item \code{purification}: $2n$ qubits; for each $i$, $H(i)$ then $\mathrm{CX}(i,i{+}n)$, i.e.\ $n$ Bell pairs
(system $i$ $\leftrightarrow$ reference $i{+}n$); order parameter is the code dimension $k=S(\text{system})$.
\item \code{single\_ref}: $n{+}1$ qubits; $H(n{-}1)$ then $\mathrm{CX}(n{-}1,n)$, i.e.\ one reference Bell-paired
with system qubit $n{-}1$; order parameter is the reference entropy $S(\{n\})\in\{0,1\}$.
\end{itemize}
The setup ops are deterministic and applied immediately (the inspector renders exactly this list). The optional
\textbf{global scramble} (Sec.~\ref{sec:runner}) then applies one maximally-scrambling random Clifford
$\code{random\_symplectic}(k)$ to the \code{scramble\_qubits()}, putting the state deep in the volume-law phase
before the monitored evolution.

% ======================================================================
\section{The execution loop}
\label{sec:runner}
\code{SimulationRunner.run(sample\_seed, *, analyses, save\_tableau, checkpoint\_layers,
checkpoint\_callback)} (\file{runner.py}) is the heart. At
construction the runner fetches the Clifford table \code{sp4\_table()} and prebuilds one gate LUT per index
(\code{\_gate\_luts\_for}, process-cached by table content). One realization proceeds:
\begin{enumerate}[itemsep=2pt]
\item \textbf{Streams \& picture.} The measurement-outcome generator
\code{sim\_rng = default\_rng([base\_seed, sample\_seed, \_MEAS\_STREAM\_KEY])} is created and handed to
\code{setup\_picture}; the \code{CircuitBuilder} gets its own independent construction stream
(Sec.~\ref{sec:records}).
\item \textbf{Global scramble} (if \code{scramble\_qubits()} is non-empty): an independent stream
\code{default\_rng([base\_seed, sample\_seed, \_SCRAMBLE\_STREAM\_KEY])} draws
\code{random\_symplectic(k)}, applied via \code{sim.apply\_clifford(\dots, qubits=scr\_qubits)}.
\item \textbf{Warmup} (gate-only layers from \code{warmup\_layers\_iter()}): each ordered layer is applied by
\code{sim.\_dispatch\_2q\_batch(gate\_pairs, cliff\_indices, lut\_array)}.
\item \textbf{Measured phase} (up to \code{total\_layers()}): for each \code{CircuitLayer}, first the
\emph{gate step} (\code{\_dispatch\_2q\_batch}), then the \emph{measurement step}
(\code{sim.\_measure\_z\_batch(meas\_qubits)}, outcomes drawn from \code{sim\_rng}). The batch kernels execute
the same scalar operations in the same order and preserve hybrid mode-switch boundaries and RNG consumption.
\end{enumerate}

\begin{lstlisting}[caption={The measured-phase inner loop (\file{runner.py}).},captionpos=b]
for layer in builder.layers():
    sim._dispatch_2q_batch(layer.gate_pairs, layer.cliff_indices,
                           lut_array)       # ordered gate step
    sim._measure_z_batch(layer.meas_qubits) # ordered measurement step
    abar_sum += sim.active_count()
    if track_op and layer.n_measurements > 0:
        ...  # recompute the order parameter (with LOCC check-skipping)
\end{lstlisting}

\paragraph{Order parameter \& early stop.} When \code{depth\_mode == until\_purified} or
\code{record\_time\_series}, the order parameter is read after any measured layer. Because it is an entanglement
entropy with a reference that is never gated, gates cannot change it and a single measurement changes it by at
most $1$ (LOCC monotonicity); the runner therefore \emph{skips} re-checks until the measurement budget since the
last check is spent, and the recorded first-zero layer $\tau_p$ stays exact. On reaching $0$ the loop breaks and
records \code{purified\_at\_layer}. It also accumulates $\bar a$ (mean active generators per qubit) as a density
diagnostic.

\paragraph{Depth checkpoints.} Requested checkpoint indices are one-based measured-layer counts and are evaluated
after that layer's gates and measurements. Without a callback, the runner stores an independent symplectic-tableau
snapshot. With \code{checkpoint\_callback}, it instead calls the read-only function as
\code{callback(sim, spec, layer)} and stores the returned value. Layers outside the executed range are ignored.
Returning the public \code{CHECKPOINT\_STOP} sentinel stops after the current layer's bookkeeping; the sentinel
itself is not stored. The final observable pass still describes the live stopped state, and
\code{purified\_at\_layer} is set by this generic stop only when the final reference order parameter is actually
zero. Sparse checkpoints therefore record the first \emph{observed} zero checkpoint, not necessarily the exact
unobserved transition layer.

% ======================================================================
\section{Gate application: the Clifford table and Numba kernels}
\label{sec:gates}
\paragraph{The Clifford set.} \code{sp4\_table()} (\file{\_clifford\_table.py}, built from
\file{core/symplectic.py}) returns a cached $(720,4,4)$ \code{uint8} array: the $|\!\Sp(4,\F_2)|=720$ two-qubit
symplectic matrices, in the basis $(X_{q_i},X_{q_j},Z_{q_i},Z_{q_j})$ with the row-vector convention
(transformed basis $=$ old $\cdot\,S$). This is the two-qubit Clifford group modulo Paulis, which is all the
phase-free tableau needs; the enumeration is verified against Stim.

\paragraph{Lookup tables.} A symplectic matrix is turned into a fast \emph{lookup table} once
(\code{\_build\_2q\_lut}, cached): a $16$-entry \code{uint8} array indexed by the packed old two-qubit Pauli bits
$(\mathit{xz}_i{\ll}2)\,|\,\mathit{xz}_j$, returning the packed new bits. The runner prebuilds one such LUT per
Clifford index up front.

\paragraph{Applying a gate.} \code{sim.\_dispatch\_2q(qi, qj, lut)} runs the two-qubit kernel
(\code{\_apply\_2q\_kernel}): it forms the deduplicated \emph{active set} of generators touching $q_i$ or $q_j$
(\code{inv[qi]} $\cup$ \code{inv[qj]}), and for each such generator packs its bits on $(q_i,q_j)$, looks up the
new bits, and updates \code{plt} plus the \code{supp}/\code{inv}/\code{inv\_x} indices where a Pauli appears or
disappears. The cost is $O(|\code{inv[qi]}|+|\code{inv[qj]}|)$ --- only the affected generators are touched, which
is the whole point of the sparse representation. \code{apply\_clifford(S, qubits)} generalizes this to an
arbitrary-arity symplectic (used for the scramble). With \code{use\_numba}, byte-identical \code{@njit} kernels
(\file{core/numba\_kernels.py}) replace the Python ones.

% ======================================================================
\section{Measurement}
\label{sec:measure}
\paragraph{Which qubits fire.} \code{sample\_measurements\_detailed(mode, n, p, gate\_pairs, rng, count)}
(\file{measurements.py}) first selects \emph{candidates} by mode --- \code{bernoulli} (all $n$ qubits),
\code{gated} (qubits in this layer's gate pairs), \code{random\_pair} (two random qubits), \code{uniform\_count}
(\code{count} random qubits) --- then flips a Bernoulli$(p)$ coin per candidate (\code{rng.random(len)}), and
returns the sorted candidates and the fired subset.

\paragraph{The tableau update.} \code{sim.measure\_z(q)} implements the Aaronson--Gottesman $Z$-measurement via
\code{\_measure\_z\_kernel}. It first checks \code{inv\_x[q]} for a \emph{stabilizer} that anticommutes with
$Z_q$: if none, the outcome is deterministic and the tableau is unchanged. Otherwise the outcome is random: a
pivot anticommuter is chosen (\code{pivot\_mode}: \code{min\_weight} keeps generators sparse, \code{first} is
cheaper), the pivot is XORed (\code{\_sparse\_xor\_rows}) into every other anticommuter, the pivot's old row is
copied into its destabilizer slot, and the pivot is replaced by the pure $Z_q$ generator. The random bit is
drawn from the simulator's independent \code{sim\_rng}.

% ======================================================================
\section{Observables and order parameters}
\label{sec:observables}
All entanglement observables reduce to a GF(2) rank (\file{core/observables.py}). \code{subsystem\_rank(sim, A)}
extracts the stabilizer block restricted to the columns of $A$ (its $x$- and $z$-bits) and computes
$\mathrm{rank}_{\F_2}$ by bit-packed Gaussian elimination (\code{gf2\_rank\_bits}, $7$--$27\times$ faster than
dense). The entanglement entropy is $S(A)=\mathrm{rank}([X|Z]|_A)-|A|$ (Fattal et al.); the
\emph{code dimension} is $k=S(\text{system})$; the \emph{reference entropy} is $S(\text{reference})$. The runner's
\code{\_order\_parameter(sim, spec)} dispatches on \code{spec.order\_parameter} to \code{code\_dimension} (purif\-
ication), \code{entanglement\_entropy(reference)} (single\_ref), or the half-cut entropy (pure\_state). The
single-reference purification time $\tau_p$ is the first layer at which this order parameter hits $0$.

% ======================================================================
\section{Records and reproducibility}
\label{sec:records}
A run returns a \code{SampleRecord} (\file{records.py}): identity (\code{sample\_seed}, \code{picture}),
diagnostics (\code{total\_layers}, \code{total\_gates}, \code{total\_measurements}, expected/actual
gate-to-measurement ratio), observables (\code{code\_dimension} / \code{ref\_entropy} / \code{entropy\_half\_cut},
whichever apply), \code{purified\_at\_layer}, \code{mean\_active\_generators} ($\bar a$), an optional
\code{time\_series}, an optional \code{final\_tableau}, any custom \code{analyses}, and either optional
\code{checkpoint\_tableaux} or callback-produced \code{checkpoint\_values} keyed by measured layer.

\paragraph{Three decoupled RNG streams.} Reproducibility schema v2 keys every per-sample generator by the ordered
pair \code{(base\_seed, sample\_seed)}, never by their collision-prone scalar sum. Additional stream components
keep measurement outcomes and scrambling independent of construction, so toggling one feature does not perturb
another:
\begin{enumerate}[itemsep=1pt]
\item \textbf{Circuit construction} --- \code{default\_rng([base\_seed, sample\_seed])} in the
\code{CircuitBuilder}, consumed in the fixed order gate-placement $\to$ Clifford-indices $\to$
measurement-candidates.
\item \textbf{Measurement outcomes} --- \code{default\_rng([base\_seed, sample\_seed, 0x6D656173])} (``meas'') on
the \code{SparseGF2} instance.
\item \textbf{Global scramble} --- \code{default\_rng([base\_seed, sample\_seed, 0x73637262])} (``scrb'').
\end{enumerate}
A stochastic graph string is realized separately from \code{base\_seed} alone and is therefore fixed across all
sample seeds in the same configuration.
A ``cell'' of a sweep is a \code{(config, sample\_seed)} pair, so results are independent of the number of
parallel workers.

% ======================================================================
\section{Compatibility and additive extensions}
\label{sec:changes}
The scientific core remains the same phase-free stabilizer simulation described above: the tableau representation,
Clifford action, measurement update, and GF(2)-rank observables do not depend on the optional scheduling and
checkpoint interfaces. The current extensions are deliberately localized:
\begin{itemize}[itemsep=2pt]
\item \textbf{Scheduling.} \code{all\_edges} adds a deterministic full-edge layer, and
\code{total\_layers\_override} supplies an exact scheduled depth. Existing gating and depth formulas retain their
documented meanings when those options are unused.
\item \textbf{Observation.} Checkpoint tableaux and read-only callbacks expose intermediate states without changing
the gate or measurement streams. \code{CHECKPOINT\_STOP} adds checkpoint-granular early stopping while preserving
final-state observables.
\item \textbf{Validation and persistence.} Configuration validation rejects ambiguous booleans, lossy integer
coercions, and nonfinite probabilities. Analysis writers use unique same-directory temporary files so overlapping
atomic writes cannot overwrite one another's payload.
\item \textbf{Reproducibility.} RNG schema v2 is an intentional stream-definition change from the former
\code{base\_seed + sample\_seed} convention. Pair-keyed streams remove seed-sum collisions; old and new payloads
must therefore carry their schema/provenance rather than being assumed bit-identical.
\end{itemize}
These interfaces are covered by parity tests across pictures and gating modes, including nonperturbation of the
executed circuit, exact-depth prefixes, checkpoint-stop bookkeeping, and equality of a final checkpoint tableau
with \code{final\_tableau}. The project changelog is the authoritative record of version-specific behavior.

\end{document}
